\section{The appointed texts}

\label{sec:appointed}

The Latin follows the 1962 Missal. Scriptural English is the Douay--Rheims (Challoner) biblical counterpart; it preserves the biblical wording rather than supplying a translation of each liturgical adaptation. The English prayers and Preface come from the Philadelphia 1861 hand missal published by Eugene Cummiskey. Their historical English is reproduced as study text. Psalm references give Vulgate numbering first and Hebrew numbering in parentheses. The precise source limits appear in the terminal appendix.

\subsection{Introit}

\label{proper-introit}

\properlocus{1962 Missal, printed 421 / PDF502, no.1759.}

\begin{latinproper}

Dicit Dominus: Ego cogito cogitationes pacis, et non afflictionis: invocabitis me, et ego exaudiam vos: et reducam captivitatem vestram de cunctis locis.

Ps. 84, 2. Benedixisti, Domine, terram tuam: avertisti captivitatem Iacob.

V. Gloria Patri.

\end{latinproper}

\begin{englishwitness}{Douay--Rheims: biblical counterpart}

11. For I know the thoughts that I think towards you, saith the Lord, thoughts of peace, and not of affliction, to give you an end and patience.

12. And you shall call upon me, and you shall go: and you shall pray to me, and I will hear you.

14. And I will be found by you, saith the Lord: and I will bring back your captivity, and I will gather you out of all nations, and from all the places to which I have driven you out, saith the Lord: and I will bring you back from the place to which I caused you to be carried away captive.

Psalm 84(85):2. Lord, thou hast blessed thy land: thou hast turned away the captivity of Jacob.

\end{englishwitness}

\subsection{Collect}

\label{proper-collect}

\properlocus{1962 Missal, printed 421 / PDF502, no.1760.}

\begin{latinproper}

Praesta, quaesumus, omnipotens Deus: ut, semper rationabilia meditantes, quae tibi sunt placita, et dictis exsequamur et factis. Per Dominum.

\end{latinproper}

\begin{englishwitness}{Cummiskey 1861: historical English}

Grant, we beseech thee, O Almighty God, that being always intent upon what is reasonable and just, we may, both in word and deed, perform what is acceptable to thee. Thro’.

\end{englishwitness}

\properlocus{Cummiskey, Sixth Sunday after Epiphany, p.~102 (PDF p.~113).}

\subsection{Epistle --- 1 Thessalonians 1:2–10}

\label{proper-epistle}

\properlocus{1962 Missal, printed 421–422 / PDF502–503, no.1761.}

\begin{latinproper}

Fratres: Gratias agimus Deo semper pro omnibus vobis, memoriam vestri facientes in orationibus nostris sine intermissione, memores operis fidei vestrae, et laboris, et caritatis, et sustinentiae spei Domini nostri Iesu Christi, ante Deum et Patrem nostrum: scientes, fratres, dilecti a Deo, electionem vestram: quia Evangelium nostrum non fuit ad vos in sermone tantum, sed et in virtute, et in Spiritu Sancto, et in plenitudine multa, sicut scitis quales fuerimus in vobis propter vos. Et vos imitatores nostri facti estis, et Domini, excipientes verbum in tribulatione multa, cum gaudio Spiritus Sancti: ita ut facti sitis forma omnibus credentibus in Macedonia, et in Achaia. A vobis enim diffamatus est sermo Domini, non solum in Macedonia, et in Achaia, sed et in omni loco fides vestra, quae est ad Deum, profecta est, ita ut non sit nobis necesse quidquam loqui. Ipsi enim de nobis annuntiant qualem introitum habuerimus ad vos: et quomodo conversi estis ad Deum a simulacris, servire Deo vivo et vero, et exspectare Filium eius de caelis (quem suscitavit ex mortuis) Iesum, qui eripuit nos ab ira ventura.

\end{latinproper}

\begin{englishwitness}{Douay--Rheims: biblical counterpart}

2. We give thanks to God always for you all: making a remembrance of you in our prayers without ceasing,

3. Being mindful of the work of your faith and labour and charity: and of the enduring of the hope of our Lord Jesus Christ before God and our Father.

4. Knowing, brethren, beloved of God, your election:

5. For our gospel hath not been unto you in word only, but in power also: and in the Holy Ghost and in much fulness, as you know what manner of men we have been among you for your sakes.

6. And you became followers of us and of the Lord: receiving the word in much tribulation, with joy of the Holy Ghost:

7. So that you were made a pattern to all that believe in Macedonia and in Achaia.

8. For from you was spread abroad the word of the Lord not only in Macedonia and in Achaia but also in every place: your faith which is towards God, is gone forth, so that we need not to speak any thing.

9. For they themselves relate of us, what manner of entering in we had unto you: and how you turned to God from idols to serve the living and true God.

10. And to wait for his Son from heaven (whom he raised up from the dead), Jesus, who hath delivered us from the wrath to come.

\end{englishwitness}

The printed citation on p.~421 is erroneous; the appointed body is 1~Thessalonians 1:2--10, as the corresponding lesson on p.~49 confirms. The biblical greeting before the thanksgiving is outside this lesson.

\subsection{Gradual --- Psalm 43(44):8–9}

\label{proper-gradual}

\properlocus{1962 Missal, printed 422 / PDF503, no.1762.}

\begin{latinproper}

Liberasti nos, Domine, ex affligentibus nos: et eos, qui nos oderunt, confudisti. V. In Deo laudabimur tota die, et in nomine tuo confitebimur in saecula.

\end{latinproper}

\begin{englishwitness}{Douay--Rheims: biblical counterpart}

8. But thou hast saved us from them that afflict us: and hast put them to shame that hate us.

9. In God shall we glory all the day long: and in thy name we will give praise for ever.

\end{englishwitness}

\subsection{Alleluia --- Psalm 129(130):1–2}

\label{proper-alleluia}

\properlocus{1962 Missal, printed 422 / PDF503, no.1763.}

\begin{latinproper}

Alleluia, alleluia. V. De profundis clamavi ad te, Domine: Domine, exaudi orationem meam. Alleluia.

\end{latinproper}

\begin{englishwitness}{Douay--Rheims: biblical counterpart}

1. Out of the depths I have cried to thee, O Lord:

2. Lord, hear my voice. Let thy ears be attentive to the voice of my supplication.

\end{englishwitness}

\subsection{Gospel --- Matthew 13:31–35}

\label{proper-gospel}

\properlocus{1962 Missal, printed 422 / PDF503, no.1764.}

\begin{latinproper}

In illo tempore: Dixit Iesus turbis parabolam hanc: Simile est regnum caelorum grano sinapis, quod accipiens homo seminavit in agro suo: quod minimum quidem est omnibus seminibus: cum autem creverit, maius est omnibus oleribus, et fit arbor, ita ut volucres caeli veniant et habitent in ramis eius. Aliam parabolam locutus est eis: Simile est regnum caelorum fermento, quod acceptum mulier abscondit in farinae satis tribus, donec fermentatum est totum. Haec omnia locutus est Iesus in parabolis ad turbas: et sine parabolis non loquebatur eis: ut impleretur quod dictum erat per Prophetam dicentem: Aperiam in parabolis os meum, eructabo abscondita a constitutione mundi.

Credo.

\end{latinproper}

\begin{englishwitness}{Douay--Rheims: biblical counterpart}

31. Another parable he proposed unto them, saying: The kingdom of heaven is like to a grain of mustard seed, which a man took and sowed in his field.

32. Which is the least indeed of all seeds; but when it is grown up, it is greater than all herbs, and becometh a tree, so that the birds of the air come, and dwell in the branches thereof.

33. Another parable he spoke to them: The kingdom of heaven is like to leaven, which a woman took and hid in three measures of meal, until the whole was leavened.

34. All these things Jesus spoke in parables to the multitudes: and without parables he did not speak to them.

35. That it might be fulfilled which was spoken by the prophet, saying: I will open my mouth in parables, I will utter things hidden from the foundation of the world.

\end{englishwitness}

\subsection{Offertory --- Psalm 129(130):1–2}

\label{proper-offertory}

\properlocus{1962 Missal, printed 422 / PDF503, no.1765.}

\begin{latinproper}

De profundis clamavi ad te, Domine: Domine, exaudi orationem meam: de profundis clamavi ad te, Domine.

\end{latinproper}

\begin{englishwitness}{Douay--Rheims: biblical counterpart}

1. Out of the depths I have cried to thee, O Lord:

2. Lord, hear my voice. Let thy ears be attentive to the voice of my supplication.

\end{englishwitness}

\subsection{Secret}

\label{proper-secret}

\properlocus{1962 Missal, printed 422 / PDF503, no.1766.}

\begin{latinproper}

Haec nos oblatio, Deus, mundet, quaesumus, et renovet, gubernet et protegat. Per Dominum.

\end{latinproper}

\begin{englishwitness}{Cummiskey 1861: historical English}

May this oblation, O God, we beseech thee, cleanse, renew, govern and protect us.

\end{englishwitness}

\properlocus{Cummiskey, Sixth Sunday after Epiphany, anomalously numbered p.~101 (PDF p.~115).}

\subsection{Preface of the Most Holy Trinity}

\label{proper-preface}

\properlocus{1962 Missal, printed 293 / PDF374, no.1082; appointment following no.1766.}

\begin{latinproper}

V. Per omnia saecula saeculorum. R. Amen.

V. Dominus vobiscum. R. Et cum spiritu tuo.

V. Sursum corda. R. Habemus ad Dominum.

V. Gratias agamus Domino Deo nostro. R. Dignum et iustum est.

Vere dignum et iustum est, aequum et salutare, nos tibi semper et ubique gratias agere: Domine, sancte Pater, omnipotens aeterne Deus: Qui cum unigenito Filio tuo, et Spiritu Sancto, unus es Deus, unus es Dominus: non in unius singularitate personae, sed in unius Trinitate substantiae. Quod enim de tua gloria, revelante te, credimus, hoc de Filio tuo, hoc de Spiritu Sancto, sine differentia discretionis sentimus. Ut in confessione verae sempiternaeque Deitatis, et in personis proprietas, et in essentia unitas, et in maiestate adoretur aequalitas. Quam laudant Angeli atque Archangeli, Cherubim quoque ac Seraphim: qui non cessant clamare cotidie, una voce dicentes:

Sanctus, Sanctus, Sanctus Dominus Deus Sabaoth. Pleni sunt caeli et terra gloria tua. Hosanna in excelsis. Benedictus qui venit in nomine Domini. Hosanna in excelsis.

\end{latinproper}

\begin{englishwitness}{Cummiskey 1861: historical English}

P. World without end. R. Amen.

P. The Lord be with you. R. And with thy spirit. P. Lift up your hearts. R. We have lifted them up to the Lord. P. Let us give thanks to the Lord our God. R. It is meet and just.

IT is truly meet and just, right and available to salvation, that we should always, and in all places, give thanks to thee, O holy Lord, Father Almighty, eternal God.

Who together with thy only begotten Son and the Holy Ghost, art one God, and one Lord not in a singularity of one Person, but in a Trinity of one substance. For what we believe of thy glory, as thou hast revealed, the same we believe of thy Son, and of the Holy Ghost, without any difference or distinction. So that in the confession of the true and eternal Deity, we adore a distinction in the Persons, an unity in the essence, and an equality in the Majesty. Whom the angels and archangels, the cherubim also and seraphim praise: and cease not daily to cry out with one voice saying, Holy, \&c.

Holy, holy, holy, Lord God of Sabaoth. Heaven and earth are full of thy glory. Hosanna in the highest. Blessed is he that cometh in the name of the Lord, Hosanna in the highest.

\end{englishwitness}

\properlocus{Cummiskey, pp.~xxvii--xxix: dialogue, common opening, Trinity clause and Sanctus.}

\subsection{Communion --- Mark 11:24}

\label{proper-communion}

\properlocus{1962 Missal, printed 422 / PDF503, no.1767.}

\begin{latinproper}

Amen dico vobis, quidquid orantes petitis, credite quia accipietis, et fiet vobis.

\end{latinproper}

\begin{englishwitness}{Douay--Rheims: biblical counterpart}

24. Therefore I say unto you, all things, whatsoever you ask when ye pray, believe that you shall receive: and they shall come unto you.

\end{englishwitness}

\subsection{Postcommunion}

\label{proper-postcommunion}

\properlocus{1962 Missal, printed 422 / PDF503, no.1768.}

\begin{latinproper}

Caelestibus, Domine, pasti deliciis: quaesumus; ut semper eadem, per quae veraciter vivimus, appetamus. Per Dominum.

\end{latinproper}

\begin{englishwitness}{Cummiskey 1861: historical English}

Being fed, O Lord, with heavenly dainties, we beseech thee, that we may always hunger after them, as they preserve our life. Thro’.

\end{englishwitness}

\properlocus{Cummiskey, Sixth Sunday after Epiphany, anomalously numbered p.~101 (PDF p.~115).}
